\documentclass[11pt,a4paper]{article}
\usepackage[a4paper,margin=25mm,headheight=15pt]{geometry}
\usepackage[T1]{fontenc}
\usepackage[utf8]{inputenc}
\IfFileExists{libertinus.sty}{\usepackage{libertinus}}{\usepackage{lmodern}}
\usepackage{microtype}
\usepackage{amsmath,amssymb,amsthm,mathtools,bm}
\usepackage{booktabs,longtable,tabularx,array}
\usepackage{xcolor,fancyhdr,needspace,xurl,listings}
\usepackage{hyperref}
\usepackage[nameinlink,noabbrev,capitalise]{cleveref}
\definecolor{ink}{RGB}{35,48,57}
\definecolor{accent}{RGB}{56,86,87}
\definecolor{lightshade}{RGB}{246,247,244}
\hypersetup{colorlinks=true,linkcolor=accent,citecolor=accent,urlcolor=accent,
 pdftitle={Real-branch inversion in power-logarithmic scales},
 pdfauthor={Technical study prepared for Vladimir Reshetnikov},
 pdfsubject={Constructive asymptotic inversion, convergence, and a Wolfram Language reference package}}
\pagestyle{fancy}\fancyhf{}
\fancyhead[L]{\small Real-branch asymptotic inversion}
\fancyhead[R]{\small Theory, algorithms, and validation}
\fancyfoot[C]{\thepage}
\renewcommand{\headrulewidth}{0.3pt}
\setlength{\parindent}{1.1em}
\setlength{\parskip}{0.25em}
\setlength{\emergencystretch}{2.5em}
\allowdisplaybreaks[2]
\numberwithin{equation}{section}
\setcounter{tocdepth}{2}
\newtheorem{theorem}{Theorem}[section]
\newtheorem{lemma}[theorem]{Lemma}
\newtheorem{proposition}[theorem]{Proposition}
\newtheorem{corollary}[theorem]{Corollary}
\theoremstyle{definition}\newtheorem{definition}[theorem]{Definition}
\newtheorem{example}[theorem]{Example}
\theoremstyle{remark}\newtheorem{remark}[theorem]{Remark}
\newcommand{\R}{\mathbb R}\newcommand{\C}{\mathbb C}\newcommand{\N}{\mathbb N}
\newcommand{\OO}{\mathcal O}\newcommand{\dd}{\mathrm d}
\newcommand{\D}{\partial_\ell}
\newcommand{\poly}{\mathsf P}\newcommand{\supp}{\operatorname{supp}}
\newcommand{\val}{\operatorname{val}}
\newcommand{\e}{\mathrm e}
\newcommand{\code}[1]{\texttt{#1}}
\lstdefinestyle{wl}{basicstyle=\ttfamily\small,columns=fullflexible,
 keepspaces=true,breaklines=true,breakatwhitespace=false,
 backgroundcolor=\color{lightshade},frame=single,rulecolor=\color{black!15},
 xleftmargin=0.6em,xrightmargin=0.6em,framesep=0.55em,
 showstringspaces=false,upquote=true,tabsize=2}
\lstset{style=wl}

\begin{document}
\begin{titlepage}
\vspace*{16mm}
{\color{accent}\large CONSTRUCTIVE ASYMPTOTIC ANALYSIS\par}
\vspace{8mm}
{\color{ink}\Huge\bfseries Real-branch inversion\par
 in power--logarithmic scales\par}
\vspace{8mm}
{\Large Explicit coefficients, rigorous remainders, convergence,\par
 and a Wolfram Language implementation\par}
\vspace{12mm}
{\large A technical study prepared for Vladimir Reshetnikov\par}
{\large 7 September 2026\par}
\vspace{13mm}
\begin{abstract}
We construct the positive-real inverse germ of
\[
 f(x)=a x^p\left(1+\sum_{i=1}^m x^{\delta_i}P_i(\log x)\right),
 \qquad a,p,\delta_i>0,
\]
where the $P_i$ are real polynomials. A normalization fixes the branch
before any expansion is attempted. An explicit multi-index
Lagrange--B\"urmann formula produces every inverse coefficient using only
polynomial differentiation. It works with irrational exponents and
correctly combines coincident exponent sums. A separate truncated
fixed-point algorithm computes the same coefficients by composition.

For this finite-input class the resulting grouped power--logarithmic
expansion is not merely formal: it converges for all sufficiently small
positive arguments. We prove this through a finite-dimensional analytic
lifting, provide a usable complex contraction criterion, and derive a
log-sensitive first-omitted-layer remainder. The two motivating examples,
$x+x^2(1+\log x)$ and $x+x^{\sqrt2}$, receive explicit formulas and detailed
branch analyses. The accompanying reference package implements exact
algebraic exponents, symbolic real coefficients under assumptions, powers
of the inverse, and formal residual checks. Its validation record
separates independently executed exact-arithmetic checks from native
Wolfram tests that were supplied but could not be run in the available
execution environment.
\end{abstract}
\vfill
{\small Deliverable version 1.0.0\quad\textbullet\quad
 Self-contained proofs and reproducible validation sources}
\end{titlepage}
\tableofcontents
\newpage

\section{The problem and the two answers}
\label{sec:answers}

The motivating question \cite{question} asks for the real inverse near the
positive endpoint $0$, rather than a Taylor expansion of an unspecified
complex inverse branch. Its first function is
\begin{equation}
 f_1(x)=x+x^2(1+\log x),\qquad x>0.
 \label{eq:f1}
\end{equation}
Its second example is
\begin{equation}
 f_2(x)=x+x^{\sqrt2},\qquad x>0.
 \label{eq:f2}
\end{equation}
All logarithms of positive real arguments in this article are real natural
logarithms. Unless another limit is stated, $y\to0^+$.

Write $g_i=f_i^{-1}$ and $L=\log y$. The first answer is
\begin{align}
 g_1(y)={}&y-y^2(1+L)+y^3(1+L)(3+2L)\nonumber\\
 &-\frac{y^4}{2}(1+L)(10L^2+31L+23)\nonumber\\
 &+\frac{y^5}{6}(1+L)(84L^3+398L^2+607L+299)
 +\OO\!\left(y^6(1+|L|)^5\right).
 \label{eq:first-answer}
\end{align}
In particular, after retaining only $y-y^2(1+\log y)$, the correct
ordinary big-$O$ bound is
\begin{equation}
 g_1(y)=y-y^2(1+\log y)
       +\OO\!\left(y^3(1+|\log y|)^2\right),
 \label{eq:log-remainder}
\end{equation}
not $\OO(y^3)$. In fact the omitted term divided by $y^3$ tends to
$+\infty$. This is an interpretation issue, not an objection to a formal
notation that explicitly declares it is grading by powers while allowing
logarithmic coefficients. The package never gives an ordinary
\code{O[y]\^{}3} object this nonstandard meaning.

Put $\alpha=\sqrt2$. The second answer, agreeing with the requested
four-term form in the archived question, is
\begin{align}
 g_2(y)={}&y-y^\alpha+\alpha y^{2\alpha-1}
 -\frac{6-\alpha}{2}y^{3\alpha-2}
 +\frac{17\alpha-12}{3}y^{4\alpha-3}\nonumber\\
 &+\frac{153\alpha-305}{12}y^{5\alpha-4}
 +\OO\!\left(y^{6\alpha-5}\right).
 \label{eq:second-answer}
\end{align}
The increment between successive exponents is $\alpha-1$, not an integer
or a rational fraction. There is no need to approximate it by a rational
number.

The rest of the article gives an all-orders construction, not just these
initial terms. The implementation covers a broader finite-input class,
including nonunit leading powers, multiple irrational increments,
logarithmic amplitudes of arbitrary finite degree, and powers of the
inverse. It is deliberately not advertised as a universal inversion
engine for every elementary function or every general-support transseries.

\subsection{What is proved and what is implemented}

The mathematical theorems allow arbitrary fixed positive real numbers
$p$ and $\delta_i$. The package restricts exponents, observable powers,
and cutoffs to \emph{exact real algebraic numbers}. This makes exponent
ordering and collisions an effective, exact operation. Coefficients may
be more general exact real constants, or real symbolic parameters with
sufficient explicit assumptions.

The construction theorem, the branch theorem, and the convergence and
remainder theorems below are proved for the stated class. Native Wolfram
execution of the package was unavailable: its remote evaluator returned
an HTTP~404 error, and no local Wolfram kernel was installed. The supplied
Python reference implementation independently passed 104 exact symbolic
checks; 11 high-precision numerical comparisons were also executed.
The package has a separate 95-test MUnit suite, but those tests are
\emph{unrun}, not counted as passed. Detailed reproduction instructions
appear in \cref{sec:validation}.

\section{Choose the branch before choosing the series}
\label{sec:branch}

\subsection{Why a complex Taylor series can answer the wrong question}

The archived question records an expansion with constant term
\[
 c=\exp\!\left(-1+W_{-1}(-\e)\right).
\]
This is not a small positive real number. The equation for a nonzero zero
of $f_1$ is $1+x(1+\log x)=0$. Setting $u=1+\log x$ gives
$u\e^u=-\e$, so suitable complex Lambert branches produce such zeros.
Expanding a local inverse around one of those regular complex zeros can
be perfectly coherent while having nothing to do with the inverse germ
that satisfies $g(y)\to0^+$.

A condition on the target such as $y>0$ does not, by itself, specify which
source germ is being inverted. The useful branch data are
\begin{equation}
 x>0,\quad y>0,\quad x\to0,\quad
 \frac{x}{(y/a)^{1/p}}\to1.
 \label{eq:branch-data}
\end{equation}
The last condition distinguishes the desired normalized germ. No
complex-root enumeration is necessary.

The behavior reported in the question is a historical observation about
the displayed commands. It is not evidence that every current Wolfram
version must behave the same way. The official documentation describes
\code{InverseSeries} and \code{AsymptoticSolve} as general tools in their
respective domains \cite{wl-inverse,wl-asymptotic}. Our purpose is to give
an explicit algorithm with a precise supported class and branch contract,
not to benchmark unexecuted current built-ins.

\subsection{A branch theorem for the entire implemented class}

\begin{theorem}[The positive inverse germ]
\label{thm:branch}
Suppose
\begin{equation}
 f(x)=a x^p(1+H(x)),\qquad
 H(x)=\sum_{i=1}^m x^{\delta_i}P_i(\log x),
 \label{eq:model}
\end{equation}
where $a,p,\delta_i>0$ and each $P_i\in\R[\ell]$. There is an
$\varepsilon>0$ such that $f$ is positive and strictly increasing on
$(0,\varepsilon)$. It has a unique inverse on its image, and this inverse
satisfies \cref{eq:branch-data}.
\end{theorem}

\begin{proof}
Every $x^\delta P(\log x)$ with $\delta>0$ tends to zero. Moreover,
\begin{equation}
 xH'(x)=\sum_i x^{\delta_i}
  \left(\delta_iP_i(\log x)+P_i'(\log x)\right)=o(1).
 \label{eq:derivative-small}
\end{equation}
Thus $1+H(x)>0$ eventually and
\[
 f'(x)=a x^{p-1}\{p(1+H(x))+xH'(x)\}
       =ap x^{p-1}(1+o(1))>0.
\]
Continuity, the limit $f(x)\to0$, and strict monotonicity give the local
inverse. From $y=a g(y)^p(1+o(1))$ and positivity, taking the real $p$th
root gives $g(y)/(y/a)^{1/p}\to1$.
\end{proof}

For \cref{eq:f1} the conclusion is global. Indeed,
\begin{equation}
 f_1'(x)=1+x(3+2\log x)\ge 1-2\e^{-5/2}>0,
 \qquad x>0.
 \label{eq:global-slope}
\end{equation}
The minimum follows by differentiating $x(3+2\log x)$. Since $f_1(0^+)=0$
and $f_1(x)\to\infty$ as $x\to\infty$, it is a bijection of $(0,\infty)$
onto itself. Likewise $f_2'(x)=1+\sqrt2\,x^{\sqrt2-1}>1$, so its positive
inverse is globally unique.

Although $g_1(0)=0$ gives a $C^1$ extension with $g_1'(0)=1$, that
extension is not $C^2$:
\[
 g_1''(y)=-\frac{5+2\log g_1(y)}{f_1'(g_1(y))^3}\longrightarrow+\infty.
\]
An ordinary Taylor series at the endpoint is therefore the wrong
representation even after the branch has been fixed.

\section{The finite power--logarithmic algebra}
\label{sec:algebra}

\subsection{Locally finite exponent supports}

Let
\begin{equation}
 \Sigma=\left\{\sum_{i=1}^m k_i\delta_i:
          \bm k=(k_1,\ldots,k_m)\in\N^m\right\},
 \qquad \delta_* =\min_i\delta_i,
 \label{eq:semigroup}
\end{equation}
where $0\in\N$. For every finite $B$, the set $\Sigma\cap[0,B]$ is finite:
$k_i\delta_i\le B$ bounds each $k_i$. This is true even if the additive
\emph{group} generated by the $\delta_i$ is dense in $\R$. Nonnegative
coefficients, rather than arbitrary integer coefficients, are essential.

Several distinct multi-indices can give the same weight. For example,
with increments $1/2$ and $1$, the weight $1$ is produced both by $(2,0)$
and by $(0,1)$. Their coefficient polynomials must be added before a term
or a remainder is reported. Algebraic equivalences such as
$\sqrt8/2=\sqrt2$ must also be recognized exactly.

A normalized formal expression has the form
\begin{equation}
 U(t)=\sum_{\Delta\in\Sigma}t^\Delta Q_\Delta(\ell),
 \qquad \ell=\log t,\quad Q_\Delta\in\R[\ell].
 \label{eq:formal-algebra}
\end{equation}
At each bounded weight there are only finitely many products to compute.
The coefficient ring is a polynomial ring in $\ell$; it is not a ring of
coefficients constant with respect to differentiation in $t$.

For genuine asymptotic comparison, a smaller power of $t$ dominates every
fixed logarithmic degree at a larger power. Within one power, the highest
nonzero logarithmic degree dominates. For a nonzero polynomial $Q$ of
degree $d$,
\[
 t^{\alpha+\eta}Q(\log t)=o(t^\alpha),\qquad \eta>0.
\]
This observation explains both the finite cutoff algorithm and the need
to preserve logarithmic degrees in ordinary remainder bounds.

\subsection{Operations used by the implementation}

Addition merges equal powers and adds their polynomials. Multiplication
is finite convolution at every cutoff:
\[
 [t^\Delta](UV)=\sum_{\beta+\gamma=\Delta}Q_\beta(\ell)R_\gamma(\ell).
\]
For a positive-order jet $U$, generalized binomial and logarithmic
expansions are well-defined by
\begin{align}
 (1+U)^r&=\sum_{n\ge0}\binom rn U^n,\label{eq:binomial}\\
 \log(1+U)&=\sum_{n\ge1}\frac{(-1)^{n+1}}nU^n.
 \label{eq:log-one-plus}
\end{align}
If $\val U\ge\delta>0$, only $n\le B/\delta$ can contribute to weights
at most $B$. This makes those infinite formulas finite algorithms.

The differentiation rule is
\begin{equation}
 \frac{\dd}{\dd t}\{t^\beta Q(\ell)\}
 =t^{\beta-1}(\beta+\D)Q(\ell).
 \label{eq:derivative-operator}
\end{equation}
Iterating it gives
\begin{equation}
 \frac{\dd^r}{\dd t^r}\{t^\beta Q(\ell)\}
 =t^{\beta-r}\prod_{j=0}^{r-1}(\beta-j+\D)Q(\ell).
 \label{eq:iterated-derivative}
\end{equation}
The factors commute because they are constant-coefficient operators in
$\ell$. Differentiation does not increase logarithmic degree.

Finally, composition into $tV$ with $V=1+U$ uses identities justified on
the positive branch:
\begin{equation}
 (tV)^\delta=t^\delta V^\delta,\qquad
 \log(tV)=\ell+\log V.
 \label{eq:normalized-composition}
\end{equation}
There is no unrestricted complex \code{PowerExpand} step. Formally,
\cref{eq:binomial,eq:log-one-plus} define the right-hand side; analytically,
$V$ is positive and close to $1$.

\section{An explicit all-orders coefficient theorem}
\label{sec:main}

The classical Lagrange--B\"urmann formula \cite{dlmf} is used here with an
\emph{auxiliary perturbation parameter}. It is not applied as an ordinary
Taylor theorem at a singular logarithmic endpoint. This is the same
phase-coordinate strategy used in the supplied companion notes
\cite{companion}, specialized to a power core and polynomial logarithmic
amplitudes.

Let
\begin{equation}
 t=(y/a)^{1/p},\qquad \ell=\log t=\frac{\log(y/a)}p.
 \label{eq:normalization}
\end{equation}
For $\bm k\in\N^m$ set
\begin{equation}
 n=|\bm k|,\quad \Delta=\bm k\cdot\bm\delta,
 \quad \bm k!=\prod_i k_i!,\quad
 P_{\bm k}(\ell)=\prod_iP_i(\ell)^{k_i}.
 \label{eq:multi-notation}
\end{equation}

\begin{theorem}[Power--logarithmic inverse coefficients]
\label{thm:coefficients}
For the positive inverse germ $g$ of \cref{eq:model}, and any fixed
$q\in\R$, one has the formal expansion
\begin{equation}
 \boxed{\displaystyle
 g(y)^q=t^q+
 \sum_{\bm k\ne\bm0}t^{q+\Delta}
 \frac{(-1)^n q}{p^n\bm k!}
 \prod_{j=1}^{n-1}(q+\Delta+pj+\D)P_{\bm k}(\ell).}
 \label{eq:master}
\end{equation}
The product is the identity for $n=1$. Equal values of $\Delta$ are
combined. At every finite power cutoff this is a finite computation.
For the exact finite model \cref{eq:model}, the expansion converges to
$g(y)^q$ for all sufficiently small $y>0$.
\end{theorem}

The proof of the coefficient identity follows immediately below.
Convergence, including the justification for using physical marker
values rather than merely small formal markers, is proved independently
in \cref{sec:convergence}. The special case $q=0$ gives the exact constant
$1$; its apparent correction terms vanish through the factor $q$.

\subsection{Derivation by a power-core coordinate}

Put $s=x^p$ and $Y=y/a$. With independent markers $\varepsilon_i$ define
\[
 h(s)=\sum_i\varepsilon_i\,
 s^{1+\delta_i/p}P_i\!\left(\frac{\log s}{p}\right).
\]
The equation becomes $Y=s+h(s)$. For fixed $Y>0$, all functions involved
are analytic on a sufficiently small complex disk about $s=Y$, using the
logarithm continued from the positive axis. For the local solution near
$s=Y$, Lagrange--B\"urmann applied to the observable $s^{q/p}$ gives
\begin{equation}
 s^{q/p}=Y^{q/p}+
 \sum_{n\ge1}\frac{(-1)^n}{n!}
 \frac{\dd^{n-1}}{\dd Y^{n-1}}
 \left\{\frac qpY^{q/p-1}h(Y)^n\right\}.
 \label{eq:LB-observable}
\end{equation}
This is an identity in small markers and also a formal marker identity.
The multinomial coefficient for $\bm\varepsilon^{\bm k}$ in $h^n$ is
$n!/\bm k!$. Before the $n-1$ derivatives, its power of $Y$ is
\[
 \frac{q+\Delta}{p}+n-1.
\]
Each derivative lowers this power by one and acts on the logarithmic
polynomial by $p^{-1}\D$. Applying \cref{eq:iterated-derivative} leaves
$Y^{(q+\Delta)/p}=t^{q+\Delta}$ and the operator factors
\[
 \prod_{j=1}^{n-1}\left(\frac{q+\Delta}{p}+j+\frac1p\D\right).
\]
Together with the initial $q/p$ and the multinomial factor, this is exactly
\cref{eq:master}. Local finiteness permits collection by the intended
power weights. The convergence argument in \cref{sec:convergence} then
allows every marker to be set to $1$ for sufficiently small $t$.

\subsection{First and second nonlinear orders}

The first inverse correction is
\begin{equation}
 g(y)=t-\frac1p\sum_i t^{1+\delta_i}P_i(\ell)+\cdots.
 \label{eq:first-general}
\end{equation}
At total marker degree two, a repeated sector $i$ contributes
\[
 \frac{t^{1+2\delta_i}}{2p^2}
 (1+2\delta_i+p+\D)P_i(\ell)^2,
\]
whereas distinct sectors $i<j$ contribute
\[
 \frac{t^{1+\delta_i+\delta_j}}{p^2}
 (1+\delta_i+\delta_j+p+\D)P_i(\ell)P_j(\ell).
\]
These are orders in the perturbation markers, not necessarily adjacent
orders in asymptotic dominance. A single large increment can occur after
many repetitions of a smaller increment. The algorithm therefore sorts
by the actual weight, not by total degree.

\subsection{Polynomial degree and coefficient computation}

If $d_i=\deg P_i$, the polynomial contributed by $\bm k$ has degree at
most $\sum_i k_i d_i$. For an individual operator $c+\D$, a coefficient
vector $(b_0,\ldots,b_d)$ is updated by
\begin{equation}
 \widetilde b_r=c b_r+(r+1)b_{r+1},\qquad b_{d+1}=0.
 \label{eq:operator-vector}
\end{equation}
Thus the potentially intimidating derivative in Lagrange inversion is
only a sequence of finite polynomial transformations. The package uses
Wolfram's polynomial \code{D} and \code{Expand} operations; a future
low-level coefficient-vector backend can implement
\cref{eq:operator-vector} without changing the theorem or public data
model.

\section{A second construction by fixed-point composition}
\label{sec:fixed}

Set $g(y)=tV(t)$. The normalized equation is
\begin{equation}
 V=\left(1+\sum_i t^{\delta_i}V^{\delta_i}
          P_i(\ell+\log V)\right)^{-1/p}.
 \label{eq:fixed-equation}
\end{equation}
Starting from $V_0=1$, define $V_{r+1}$ by the right-hand side, truncating
all jet operations to the requested finite weight bound.

\begin{proposition}[Finite stabilization]
\label{prop:fixed}
If two normalized jets $V$ and $W$ agree below relative weight $\beta>0$,
their images under \cref{eq:fixed-equation} agree below weight
$\beta+\delta_*$. Consequently $V_r$ agrees with the formal solution at
all weights strictly below $(r+1)\delta_*$. In particular,
$\lfloor B/\delta_*\rfloor+1$ iterations suffice for all weights at most
$B$.
\end{proposition}

\begin{proof}
For $V,W\in1+\mathfrak m$, the difference between $V^\delta$ and
$W^\delta$ is $(V-W)$ times a formal series of nonnegative order. The same
is true of $\log V-\log W$, by the binomial and logarithmic identities.
Substituting these into a polynomial $P$ preserves that property. Each
correction is then multiplied by $t^{\delta_i}$, so its difference gains
at least $\delta_*$. The outer power $-1/p$ has a nonzero constant base
and does not reduce order. Iteration gives the assertion, beginning with
$V-1$ of order at least $\delta_*$.
\end{proof}

This also proves formal uniqueness. If two solutions first differed at
weight $\beta$, applying the equation would force their difference to
start no earlier than $\beta+\delta_*$, a contradiction. Thus the
fixed-point construction and \cref{eq:master} compute the same formal
object.

For the near-identity case $p=a=1$, an alternative unnormalized iteration
is $G_{r+1}(y)=y-h(G_r(y))$, with $h=f-x$. The normalized iteration is more
uniform: it also handles $p\ne1$ and avoids taking a noninteger power of
a series with zero constant term. Its internal powers are always powers
of $1+$ a positive-order jet.

The two algorithms have different failure modes in an implementation.
A missing derivative in \cref{eq:master} can be detected by composition;
a missing cross term in composition can be detected by the explicit
multi-index formula. Their agreement is therefore a useful test, although
it does not replace the proofs.

\section{Why the grouped expansions actually converge}
\label{sec:convergence}

\subsection{A finite-dimensional analytic lifting}

Write $P_i(\ell)=\sum_{k=0}^{d_i}c_{ik}\ell^k$. Expanding
$P_i(\ell+\log V)$ gives
\begin{equation}
 t^{\delta_i}P_i(\ell+\log V)
 =\sum_{j=0}^{d_i}z_{ij}A_{ij}(\log V),\quad
 z_{ij}=t^{\delta_i}\ell^j,
 \label{eq:lift}
\end{equation}
where
\[
 A_{ij}(u)=\sum_{k=j}^{d_i}c_{ik}\binom kj u^{k-j}.
\]
There are finitely many independent variables $z_{ij}$. Consider
\begin{equation}
 F(V,\bm z)=V^p\left(1+
   \sum_{i,j}z_{ij}V^{\delta_i}A_{ij}(\log V)\right)-1.
 \label{eq:analytic-lift}
\end{equation}
Choose the complex logarithm analytic on a disk centered at $1$ that does
not meet $0$. Then $F$ is analytic near $(V,\bm z)=(1,\bm0)$,
$F(1,\bm0)=0$, and $F_V(1,\bm0)=p\ne0$.

\begin{theorem}[Convergent analytic lifting]
\label{thm:analytic}
There is an analytic function $V=\mathcal V(\bm z)$ near $\bm z=0$ with
$\mathcal V(0)=1$ satisfying \cref{eq:analytic-lift}. Its Taylor series
converges absolutely in a sufficiently small polydisk. Along the real
curve $z_{ij}=t^{\delta_i}(\log t)^j$, all variables tend to zero, so
$t\mathcal V(\bm z(t))$ is the positive inverse for sufficiently small
$t>0$. Grouping the Taylor series by exponent weights gives
\cref{eq:master}, and the grouped expansion converges absolutely after
its polynomial amplitudes are evaluated at $\ell=\log t$.
\end{theorem}

\begin{proof}
The analytic implicit function theorem gives $\mathcal V$. Each
$z_{ij}(t)$ tends to zero because $\delta_i>0$. For real sufficiently
small $t$, the solution remains real and close to $1$, hence positive;
\cref{thm:branch} identifies it with the required inverse. Every monomial
in the Taylor series has weight $\sum_i k_i\delta_i$ after substitution,
with a finite logarithmic polynomial obtained by collecting the choices
of $j$. Absolute convergence allows regrouping. Formal uniqueness from
\cref{prop:fixed}, or the marker derivation above, identifies the grouped
coefficients with \cref{eq:master}.
\end{proof}

It is useful to distinguish this theorem from a general assertion about
transseries. General transseries need not converge as ordinary sums, and
formal identities do not automatically justify evaluation of arbitrary
marker values. Our conclusion uses the \emph{finite} number of
power--logarithmic inputs and the analytic function near $V=1$. The
broader grid-based viewpoint is discussed in \cite{edgar-beginners,edgar-grids};
no unproved general-support closure statement is needed here.

\subsection{An explicit sufficient convergence criterion}

The preceding proof is qualitative. A contraction estimate gives a
computable sufficient condition at a specified positive $t$. Fix
$0<r<1$ and define
\[
 M_r=-\log(1-r),\qquad
 P_i^\#(u)=\sum_k |c_{ik}|u^k.
\]
For $\ell=\log t$ set
\begin{align}
 A(t)&=\sum_i t^{\delta_i}(1+r)^{\delta_i}
          P_i^\#(|\ell|+M_r),\label{eq:A}\\
 B(t)&=\frac1{1-r}\sum_i t^{\delta_i}(1+r)^{\delta_i}
 \left\{\delta_iP_i^\#(|\ell|+M_r)
                  +(P_i^\#)'(|\ell|+M_r)\right\}.
 \label{eq:B}
\end{align}
On $|V-1|\le r$, these bound respectively $|H(tV)|$ and its derivative
with respect to $V$. The estimates follow from
$|\log V|\le M_r$, $|V|\le1+r$, and $|V|\ge1-r$.

\begin{proposition}[Marker disk and tail bound]
\label{prop:contraction}
Let $R>1$. If $RA(t)<1$ and
\begin{align}
 \frac{RA(t)}p(1-RA(t))^{-1/p-1}&\le r,\label{eq:disk-test}\\
 \kappa:=\frac{RB(t)}p(1-RA(t))^{-1/p-1}&<1,\label{eq:contraction-test}
\end{align}
then the normalized inverse is analytic in the common perturbation marker
$\varepsilon$ for $|\varepsilon|\le R$ and remains in $|V-1|\le r$.
If all multi-index terms of total degree at most $N$ are retained, then
\begin{equation}
 |g(y)-G_N(y)|\le
 t\,\frac{rR^{-(N+1)}}{1-R^{-1}}.
 \label{eq:geometric-tail}
\end{equation}
This bounds a total-marker-degree truncation, not an arbitrary weighted
cutoff without further accounting.
\end{proposition}

\begin{proof}
The map $T(V)=(1+\varepsilon H(tV))^{-1/p}$ is analytic on the disk because
$|\varepsilon H|\le RA<1$. Integrating its derivative with respect to
$\varepsilon H$ along the straight segment from $0$ shows that
$|T(V)-1|$ is bounded by the left side of \cref{eq:disk-test}. Its
$V$-derivative is bounded by \cref{eq:contraction-test}. The map is a
uniform contraction of the closed disk into itself. Iteration gives its
unique fixed point, analytically in $\varepsilon$. Cauchy's estimate on
$V-1$, whose absolute value is at most $r$, bounds its $n$th marker
coefficient by $rR^{-n}$. Summing the geometric tail and multiplying by
$t$ proves \cref{eq:geometric-tail}.
\end{proof}

All quantities in \cref{eq:A,eq:B} tend to zero as $t\to0^+$, so the
criterion succeeds eventually for any fixed $r$ and $R$. It is a
sufficient test; failure does not imply divergence.

\section{Power cutoffs and rigorous logarithmic remainders}
\label{sec:remainders}

The public cutoff $b$ means: retain every term whose \emph{target-variable
power} is strictly less than $b$. For an observable $g^q$, a normalized
term $t^{q+\Delta}Q_\Delta(\ell)$ has target power $(q+\Delta)/p$.
Thus the relative weight threshold is
\begin{equation}
 B=pb-q>0.
 \label{eq:relative-cutoff}
\end{equation}
Define the first possible omitted weight
\begin{equation}
 \lambda=\min\{\Delta\in\Sigma:\Delta\ge B\}.
 \label{eq:frontier}
\end{equation}
The associated coefficient is the \emph{combined} polynomial
$Q_\lambda$, which can vanish through cancellation.

\begin{theorem}[First omitted layer]
\label{thm:remainder}
Let $G_{<b}$ contain every term of $g(y)^q$ of target power less than $b$.
For the exact finite model,
\begin{equation}
 g(y)^q-G_{<b}(y)
 =t^{q+\lambda}Q_\lambda(\ell)+o(t^{q+\lambda}).
 \label{eq:first-omitted}
\end{equation}
If $Q_\lambda\ne0$ and $d=\deg Q_\lambda$, then
\begin{equation}
 g(y)^q-G_{<b}(y)
 =\OO\!\left(t^{q+\lambda}(1+|\log t|)^d\right).
 \label{eq:remainder-general}
\end{equation}
If $Q_\lambda=0$, the bound $\OO(t^{q+\lambda})$ is valid, though generally
not sharp.
\end{theorem}

\begin{proof}
Local finiteness gives a positive gap between $\lambda$ and the next
possible weight. After removing finitely many layers, the remaining
Taylor series of \cref{eq:analytic-lift} has arbitrarily high total degree
in the finite variables $z_{ij}$. If $d_*=\max_i d_i$, their largest
absolute value is bounded by a constant times
$t^{\delta_*}(1+|\log t|)^{d_*}$. Choose the removed Taylor degree $N$ so
large that $(N+1)\delta_*>\lambda$. The analytic Taylor tail is then
$o(t^\lambda)$. The finitely many remaining layers above $\lambda$ are
also $o(t^\lambda)$ by their positive power gaps. Multiplication by
$t^q$ proves \cref{eq:first-omitted}. The polynomial degree gives the
stated ordinary big-$O$ bounds.
\end{proof}

\subsection{A finite way to find the omitted layer}

Let $n_*=\lceil B/\delta_*\rceil$. The semigroup contains
$n_*\delta_*\ge B$, so
\[
 \lambda\le n_*\delta_*.
\]
It therefore suffices to enumerate multi-indices satisfying
\begin{equation}
 \bm k\cdot\bm\delta\le n_*\delta_*.
 \label{eq:enumeration-bound}
\end{equation}
This gives all retained weights and the full first omitted bucket. It
also gives $|\bm k|\le n_*$. Importantly, a numerical approximation to
$\delta_*$ is neither needed nor safe at a boundary.

The package returns both the polynomial $Q_\lambda$ and its degree. If
that polynomial cancels, it does not silently claim to have found the
first \emph{nonzero} omitted layer. A zero frontier still supplies the
conservative bound in \cref{thm:remainder}. Searching farther would
improve sharpness, not the correctness of the retained expansion.

\subsection{Why a custom remainder object is used}

\code{PowerLogRemainder[t,r,d]} means only
$\OO(t^r(1+|\log t|)^d)$ on the positive axis. It is an inert descriptor,
not a \code{SeriesData} object and not an implemented algebra of error
terms. The returned association separately stores the finite expression,
its terms, the target power of the remainder, and its logarithmic degree.
Adding two such descriptors does not magically combine their bounds.
This deliberately small interface prevents an attractive display from
making an unsupported mathematical promise.

\section{The logarithmic example in detail}
\label{sec:log-example}

For $f_1$, the model data are $a=p=1$, one increment $\delta=1$, and
$P(\ell)=1+\ell$. The all-orders formula reduces to
\begin{equation}
 g_1(y)=y+\sum_{n\ge1}y^{n+1}A_n(L),\qquad
 A_n(L)=\frac{(-1)^n}{n!}
 \prod_{j=1}^{n-1}(n+1+j+\partial_L)(1+L)^n.
 \label{eq:log-all}
\end{equation}
The first five polynomials are
\begin{align*}
 A_1={}&-(L+1),\\
 A_2={}&(L+1)(2L+3),\\
 A_3={}&-\tfrac12(L+1)(10L^2+31L+23),\\
 A_4={}&\tfrac16(L+1)(84L^3+398L^2+607L+299),\\
 A_5={}&-\tfrac1{12}(L+1)
   (504L^4+3223L^3+7507L^2+7565L+2789).
\end{align*}
Every polynomial $A_n$ has the factor $L+1$: at most $n-1$
differentiations are applied to $(L+1)^n$. This is consistent with the
exact fixed point $f_1(\e^{-1})=\e^{-1}$, although the asymptotic
construction concerns the endpoint $0$ rather than that fixed point.

\subsection{The highest logarithms are Catalan coefficients}

Only the constant part of each operator contributes to the coefficient
of $L^n$. Consequently
\begin{equation}
 [L^n]A_n(L)=\frac{(-1)^n}{n!}
   \prod_{j=1}^{n-1}(n+1+j)
 =(-1)^n\frac1{n+1}\binom{2n}{n}.
 \label{eq:catalan}
\end{equation}
Thus the leading-logarithm diagonal is controlled by Catalan numbers.
This is not an empirical pattern inferred from a few coefficients; it
follows for every $n$ from the differential-operator formula.

Keeping through $A_N$ gives
\begin{equation}
 g_1(y)=y+\sum_{n=1}^N y^{n+1}A_n(\log y)
 +\OO\!\left(y^{N+2}(1+|\log y|)^{N+1}\right).
 \label{eq:log-all-remainder}
\end{equation}
The degree $N+1$ is sharp for this ordinary power--log envelope, because
its highest coefficient \cref{eq:catalan} is nonzero.

\subsection{An explicit convergence interval}

Apply \cref{prop:contraction} with $r=1/2$ and $R=2$. Here $t=y$ and
\[
 A(t)=\tfrac32t(1+|\log t|+\log2),\qquad
 B(t)=3t(2+|\log t|+\log2).
\]
Both displayed functions increase on $(0,1/100]$. Using $\log100<5$ and
$\log2<1$ gives throughout that interval
\[
 A\le\frac{21}{200},\qquad B\le\frac6{25}.
\]
The disk displacement bound and contraction constant therefore satisfy
\[
 \frac{2A}{(1-2A)^2}\le\frac{2100}{6241}<\frac12,
 \qquad
 \frac{2B}{(1-2A)^2}\le\frac{4800}{6241}<1.
\]
This proves convergence of \cref{eq:log-all} for every
$0<y\le1/100$, with the explicit, deliberately conservative bound
\begin{equation}
 \left|g_1(y)-y-\sum_{n=1}^N y^{n+1}A_n(\log y)\right|
 \le y\,2^{-(N+1)}.
 \label{eq:log-convergence-bound}
\end{equation}
The power--log bound \cref{eq:log-all-remainder} is much more informative
as $y\to0$ at fixed $N$; \cref{eq:log-convergence-bound} instead guarantees
convergence as $N\to\infty$ on an explicit interval. These are different
uses of a remainder estimate.

\subsection{A useful partial resummation}

The analytic lifting for this example has only two variables,
\[
 z=y\log y,\qquad w=y.
\]
Its equation is
\begin{equation}
 V+(z+w)V^2+wV^2\log V=1.
 \label{eq:two-variable-log}
\end{equation}
At $w=0$, the branch near $V=1$ is
\begin{equation}
 V_0(z)=\frac{2}{1+\sqrt{1+4z}}.
 \label{eq:leading-log-sum}
\end{equation}
The rationalized form avoids subtracting nearly equal numbers at $z=0$.
Differentiating \cref{eq:two-variable-log} with respect to $w$ at $w=0$
gives
\[
 \left.\frac{\partial V}{\partial w}\right|_{w=0}
 =-\frac{V_0(z)^2(1+\log V_0(z))}{1+2zV_0(z)}
 =-\frac{V_0(z)^2(1+\log V_0(z))}{\sqrt{1+4z}}.
\]
Hence, as $y\to0^+$,
\begin{equation}
 g_1(y)=yV_0(y\log y)
 -y^2\frac{V_0(y\log y)^2\{1+\log V_0(y\log y)\}}
               {\sqrt{1+4y\log y}}
 +\OO(y^3).
 \label{eq:resummed-log}
\end{equation}
Here an ordinary $\OO(y^3)$ really is valid: infinitely many leading and
subleading logarithms have been retained in exact analytic coefficient
functions. The analytic Taylor remainder in $w$ is uniform for $z$ in a
small fixed disk, and $z=y\log y$ eventually lies there. This resummation
is a theoretical and numerical alternative, not an extra package backend.
The package returns finite power--logarithmic truncations.

\section{The irrational-power example in detail}
\label{sec:irrational}

Consider the more general family
\begin{equation}
 f_\alpha(x)=x+c x^\alpha,\qquad \alpha>1,
 \label{eq:pure-family}
\end{equation}
with fixed real $c$; positivity and monotonicity hold near zero even when
$c<0$. Put $\delta=\alpha-1$. The coefficient theorem gives
\begin{equation}
 g_\alpha(y)=y\sum_{n\ge0}C_n(\alpha)(c y^\delta)^n,
 \label{eq:pure-series}
\end{equation}
where $C_0=1$ and
\begin{align}
 C_n(\alpha)
 &=\frac{(-1)^n}{n!}\prod_{j=0}^{n-2}(n\alpha-j)\label{eq:pure-product}\\
 &=\frac{(-1)^n}{1+n(\alpha-1)}\binom{n\alpha}{n}\label{eq:pure-binomial}\\
 &=(-1)^n\frac{\Gamma(n\alpha+1)}
                 {\Gamma(n+1)\Gamma(n(\alpha-1)+2)}.
 \label{eq:pure-gamma}
\end{align}
The finite product is the preferred exact computational formula. It
avoids unnecessary gamma functions and their associated simplification
and pole-cancellation issues.

The first coefficients are
\begin{align*}
 C_1&=-1,\qquad C_2=\alpha,\\
 C_3&=-\frac{\alpha(3\alpha-1)}2,\\
 C_4&=\frac{\alpha(4\alpha-1)(2\alpha-1)}3.
\end{align*}
Substituting $\alpha=\sqrt2$, $c=1$ gives \cref{eq:second-answer}. An
ordinary power cutoff $b>1$ retains exactly those integers $n\ge0$ with
$1+n\delta<b$. This is a strict inequality; using a floating-point
\code{Floor} at an exact boundary can retain or omit the wrong term.

\subsection{Convergence radius and the complex obstruction}

Write $z=c y^\delta$ and $g=yV(z)$. Then
\begin{equation}
 V+zV^\alpha=1,\qquad V(0)=1.
 \label{eq:pure-implicit}
\end{equation}
This defines an ordinary analytic function of $z$ near zero, even when
$\alpha$ is irrational. Stirling's formula applied to
\cref{eq:pure-gamma} yields
\begin{equation}
 |C_n(\alpha)|\sim
 \frac{\sqrt\alpha}{\sqrt{2\pi}(\alpha-1)^{3/2}}
 n^{-3/2}
 \left(\frac{\alpha^\alpha}{(\alpha-1)^{\alpha-1}}\right)^n.
 \label{eq:coefficient-asymptotic}
\end{equation}
Therefore the exact radius of convergence in $z$ is
\begin{equation}
 R_\alpha=\frac{(\alpha-1)^{\alpha-1}}{\alpha^\alpha}.
 \label{eq:pure-radius}
\end{equation}
The same value is suggested by the critical point of the implicit
equation: simultaneously solving $V+zV^\alpha=1$ and
$1+\alpha zV^{\alpha-1}=0$ gives
\[
 V_c=\frac\alpha{\alpha-1},\qquad z_c=-R_\alpha.
\]
The coefficient asymptotic, rather than an unproved assertion that this
is the nearest singularity, establishes the radius. Because the boundary
coefficients have size proportional to $n^{-3/2}$, the series converges
absolutely on $|z|=R_\alpha$ as well.

For $\alpha=\sqrt2$ the numerical values are
\begin{equation}
 R_\alpha\approx0.42519560665184081380,\qquad
 R_\alpha^{1/(\alpha-1)}\approx0.12686270512330968140.
 \label{eq:numeric-radius}
\end{equation}
For $c=1$, the positive inverse exists for every $y>0$, but this particular
series centered at $z=0$ has a finite convergence radius. Global
invertibility and convergence of one local expansion are separate facts.

\subsection{Why rational approximation changes the problem}

Replacing $\sqrt2$ by a nearby rational number makes a Puiseux grid
available, but it computes the inverse of a different function. If two
exponents differ by a fixed nonzero $\eta$, the ratio of their monomials
is $y^\eta$, which does not tend to $1$ as $y\to0$. At sufficiently small
$y$, even a very small fixed exponent error changes asymptotic dominance.
The package consequently rejects inexact exponents instead of silently
rationalizing them.

\section{Mixed sectors, nonlinear cores, and inverse observables}
\label{sec:examples}

\subsection{A resonant mixed example}

Let
\[
 f(x)=x+x^{3/2}+2x^2\log x.
\]
The increments are $1/2$ and $1$. At relative weight $1$, the quadratic
interaction of the first sector contributes $3/2$, while the second
sector contributes $-2\log y$. At relative weight $3/2$, a cubic term
and a mixed term must also be combined. The result is
\begin{align}
 g(y)={}&y-y^{3/2}+y^2\left(\frac32-2\log y\right)
 +y^{5/2}\left(7\log y-\frac58\right)\nonumber\\
 &+\OO\!\left(y^3(1+|\log y|)^2\right).
 \label{eq:mixed}
\end{align}
Enumerating one term for each input sector would miss both interactions.
Treating equal weights as distinct ordered monomials would also give the
wrong remainder semantics.

\subsection{A nonlinear leading power}

Take
\[
 f(x)=2x^3\left(1+x^{1/2}(1+\log x)\right),\qquad
 t=(y/2)^{1/3},\quad\ell=\log t.
\]
The positive inverse is
\begin{equation}
 g(y)=t-\frac{1+\ell}{3}t^{3/2}
       +\frac{(1+\ell)(5\ell+7)}{18}t^2
       +\OO\!\left(t^{5/2}(1+|\ell|)^3\right).
 \label{eq:nonlinear-core}
\end{equation}
The target powers are $1/3$, $1/2$, and $2/3$; the next is $5/6$.
The logarithmic coefficient is evaluated at $\log(y/2)/3$, not at
$\log y$. Omitting the constant $a$ or the factor $p$ in the logarithm
changes coefficients, even though the leading order may look plausible.

\subsection{Powers of the inverse without repeated composition}

Setting $q=2$ in \cref{eq:master} directly computes $g^2$; it does not
first construct a longer expansion of $g$ and square it. For instance,
for $f(x)=x+x^2$,
\[
 g(y)^2=y^2-2y^3+5y^4-14y^5+\OO(y^6).
\]
Negative and fractional $q$ are equally valid on the positive branch.
The only change to cutoff bookkeeping is $B=pb-q$. This makes the
observable formula useful for quantities that depend on the inverse
without requiring the inverse as an intermediate output.

\subsection{A cancelled omitted layer}

For $f(x)=x+x^2+2x^3$, the inverse starts
\[
 g(y)=y-y^2+0\,y^3+5y^4+\cdots.
\]
A strict cutoff $b=3$ retains $y-y^2$. The first possible omitted weight
is $2$, but its coefficient is zero. Reporting $\OO(y^3)$ is correct and
conservative; claiming that the error is asymptotic to a nonzero multiple
of $y^3$ would be false. The result association explicitly exposes this
cancellation in \code{"FrontierPolynomial"}.

\section{Residuals, inverse-error bounds, and forward jets}
\label{sec:residual}

\subsection{The residual-to-error principle}

\begin{proposition}[A posteriori inverse bound]
\label{prop:residual}
Let $f$ be continuously differentiable and increasing on an interval
$I$, with $f'(x)\ge m>0$ there. Suppose the true inverse value $g(y)$ and
an approximation $G(y)$ both belong to $I$. Then
\begin{equation}
 |G(y)-g(y)|\le\frac{|f(G(y))-y|}{m}.
 \label{eq:residual-bound}
\end{equation}
\end{proposition}
\begin{proof}
The mean value theorem gives $f(G)-f(g)=f'(\xi)(G-g)$ for a point between
$G$ and $g$. Taking absolute values proves the assertion.
\end{proof}

For a general power core, the relevant slope is not a fixed constant:
it is asymptotic to $ap t^{p-1}$. On a comparability interval such as
$[t/2,2t]$, sufficiently small $t$ gives a lower bound
$c t^{p-1}$ with $c>0$. Thus
\begin{equation}
 f(G)-y=\OO(t^{p+\beta}(1+|\log t|)^d)
 \quad\Longrightarrow\quad
 G-g=\OO(t^{1+\beta}(1+|\log t|)^d).
 \label{eq:slope-transfer}
\end{equation}
Forgetting the factor $t^{p-1}$ produces a wrong inverse remainder when
$p\ne1$.

For the two motivating examples, there are particularly simple global
bounds for every positive approximation $G$:
\begin{align}
 |G-g_1(y)|&\le
 \frac{|f_1(G)-y|}{1-2\e^{-5/2}},\label{eq:log-residual-bound}\\
 |G-g_2(y)|&\le |f_2(G)-y|.
 \label{eq:pure-residual-bound}
\end{align}
These are exact mathematical inequalities. Turning a floating-point
residual evaluation into a certified numerical error bar additionally
requires enclosing its rounding error, for example with validated
interval arithmetic. The included high-precision numerical experiments
are not claimed to be such interval certificates.

\subsection{What the package's formal residual checks}

For $G=tV$, the normalized residual is
\begin{equation}
 \mathcal R(t)=\frac{f(G)}{a t^p}-1
 =V^p\left(1+\sum_i t^{\delta_i}V^{\delta_i}
                 P_i(\ell+\log V)\right)-1.
 \label{eq:normalized-residual}
\end{equation}
\code{RealInverseResidual} evaluates this identity in the truncated jet
algebra. A zero result below relative weight $B$ establishes algebraic
cancellation of all checked coefficients. It is not a claim that the
exact residual is identically zero or that a numerical enclosure was
computed.

There is also a useful sign-and-slope check. If the first missing inverse
layer is $t^{1+\lambda}Q_\lambda(\ell)$, then the leading normalized
residual of the retained approximation is
\begin{equation}
 \mathcal R(t)=-p\,t^\lambda Q_\lambda(\ell)+o(t^\lambda).
 \label{eq:frontier-residual}
\end{equation}
This tests the derivative normalization, the sign of reversion, and the
frontier polynomial at once. It is among the independent exact checks.

\subsection{Using a finite forward asymptotic jet}

The implementation treats its finite input as an \emph{exact function}.
Suppose instead that a true forward function is
\[
 F(x)=a x^p(1+H_0(x)+R(x)),
\]
where $H_0$ is the finite model supplied to the package and, for some
$\rho>0$ and nonnegative integer $s$,
\begin{equation}
 R(x)=\OO(x^\rho(1+|\log x|)^s),\qquad
 xR'(x)=\OO(x^\rho(1+|\log x|)^s).
 \label{eq:forward-remainder}
\end{equation}
These assumptions imply the same positive slope asymptotic and the
existence of the inverse germ. If $g_0$ is the inverse of the finite
model, evaluating $F(g_0)-y$ and using \cref{eq:slope-transfer} gives
\begin{equation}
 F^{-1}(y)-g_0(y)
 =\OO(t^{1+\rho}(1+|\log t|)^s).
 \label{eq:jet-stability}
\end{equation}
More generally their $q$th powers differ by
$\OO(t^{q+\rho}(1+|\log t|)^s)$.

Therefore every inverse coefficient at relative weight strictly below
$\rho$ is determined by the supplied forward jet. At the boundary,
unknown forward information can alter the omitted polynomial. Simply
applying \code{Normal} to a \code{SeriesData} input and discarding its
remainder does not justify arbitrarily high inverse orders. The package
rejects an unparsed \code{SeriesData} or explicit error term; management
of a forward remainder belongs to the caller unless the exact finite
model is genuinely the intended function.

\subsection{Newton refinement}

Once a positive approximation $G$ is available, a Newton correction is
\[
 G_{\rm new}=G-\frac{f(G)-y}{f'(G)}.
\]
For the finite model, on a sufficiently small comparability interval,
$f''/f'=\OO(1/t)$. If
$G-g=\OO(t^{1+\beta}(1+|\log t|)^d)$, Taylor's theorem gives
\[
 G_{\rm new}-g=\OO(t^{1+2\beta}(1+|\log t|)^{2d}).
\]
This is an optional numerical refinement or a possible future formal
backend. It is not a third implemented package method. A numerical
implementation should keep a positive bracket when needed, rather than
assume that an unconstrained root finder will preserve the branch.

\section{The Wolfram Language package}
\label{sec:package}

\subsection{Loading and immediate use}

The package is self-contained and has no third-party Wolfram dependencies.
After extracting the archive, load its kernel file:
\begin{lstlisting}
Get["path/to/RealInverseAsymptotics/Kernel/RealInverseAsymptotics.wl"];
Clear[x, y];

r1 = RealInverseSeries[x + x^2 (1 + Log[x]), {x, 0}, y, 6];
r1["Expression"]
r1["Remainder"]

r2 = RealInverseSeries[x + x^Sqrt[2], {x, 0}, y,
                      4 Sqrt[2] - 3];
r2["Expression"]
\end{lstlisting}
For \code{r1}, target powers strictly below $6$ are retained, so the
expression is \cref{eq:first-answer} without its remainder. Its remainder
metadata has target power $6$ and logarithmic degree $5$. For \code{r2},
the cutoff is exactly the next power after the four requested terms;
that next term is not retained.

The direct model interface avoids relying on symbolic parsing:
\begin{lstlisting}
Clear[ell, y];
r = InversePowerLog[2, 3, {{1/2, 1 + ell}}, ell, y, 5/6];
r["Expression"]
\end{lstlisting}
This produces \cref{eq:nonlinear-core} through its first three terms.
The log placeholder \code{ell} is independent of the target \code{y}.
Variables used as source, target, and placeholder should have no assigned
values.

\subsection{Public functions}

\begin{longtable}{@{}>{\raggedright\arraybackslash}p{0.41\textwidth}>{\raggedright\arraybackslash}p{0.55\textwidth}@{}}
\toprule
Function & Contract \\
\midrule
\endhead
\code{PowerLogNormalForm}\newline\code{[f,x]} & Parse an exact finite power--logarithmic
expression into coefficient $a$, leading power $p$, and correction pairs
$\{\delta_i,P_i\}$. \\
\addlinespace
\code{RealInverseSeries}\newline\code{[f,\{x,0\},y,b]} & Return a branch-aware inverse
result association, keeping target powers strictly below $b$. \\
\addlinespace
\code{InversePowerLog}\newline\code{[a,p,corr,L,y,b]} & Use explicitly supplied normalized
model data; \code{corr} is a list of pairs. \\
\addlinespace
\code{RealInverseResidual}\newline\code{[result,B]} & Compute the finite normalized
residual jet at relative $t$-powers strictly below $B$. Requires the
ordinary inverse, not another observable power. \\
\addlinespace
\code{PurePowerInverseCoefficient}\newline\code{[alpha,n]} & Return
\cref{eq:pure-product} for exact real algebraic $\alpha>1$ and integer
$n\ge0$. \\
\addlinespace
\code{PowerLogRemainder}\newline\code{[t,r,d]} & Inert ordinary big-$O$ metadata,
$\OO(t^r(1+|\log t|)^d)$. \\
\bottomrule
\end{longtable}

\subsection{Result fields and their coordinate systems}

The following fields are important when using a result programmatically.
\begin{longtable}{@{}>{\raggedright\arraybackslash}p{0.34\textwidth}>{\raggedright\arraybackslash}p{0.62\textwidth}@{}}
\toprule
Field & Meaning \\
\midrule
\endhead
\code{"Expression"} & The explicit finite expression in $y$, with
logarithms evaluated at $\log(y/a)/p$. \\
\code{"Terms"} & Pairs $\{r,Q(\ell)\}$ meaning $t^rQ(\ell)$, not
$y^rQ(\log y)$. The basis is recorded in \code{"TermBasis"}. \\
\code{"Coefficient"}, \code{"LeadingPower"} & The forward leading data
$a$ and $p$. \\
\code{"Corrections"}, \code{"LogVariable"} & The normalized input and
its polynomial placeholder. \\
\code{"Cutoff"} & The strict target-variable power cutoff $b$. \\
\code{"ObservablePower"} & The $q$ in $g(y)^q$; default $1$. \\
\code{"FrontierWeight"} & The relative normalized weight $\lambda$. \\
\code{"FrontierPolynomial"} & The combined $Q_\lambda$; it may be zero. \\
\code{"FrontierExpression"} & The first possible omitted layer expressed
in $y$, supplied for nonexact results. \\
\code{"RemainderPower"} & The target-variable power $(q+\lambda)/p$. \\
\code{"RemainderLogDegree"} & The degree of $Q_\lambda$, or zero for a
cancelled frontier. \\
\code{"Remainder"} & The inert descriptor in the normalized $t$ coordinate. \\
\code{"IsExact"} & True for the exact monomial inverse or the constant
observable $q=0$; otherwise exactness is not claimed. \\
\code{"Branch"}, \code{"Assumptions"} & The positive germ and the
parameter/target conditions under which the result is intended. \\
\code{"EnumeratedMultiIndices"} & Number of simplex points enumerated,
including the zero multi-index. \\
\bottomrule
\end{longtable}

The distinction among $b$, $B$, and $\lambda$ is intentional. For
$a=p=q=1$, the simple relation $B=b-1$ can hide a coordinate mistake that
becomes visible only for a nonlinear leading power. The residual function
uses a relative normalized cutoff; the inverse function uses an absolute
target-power cutoff.

\subsection{Options and symbolic parameters}

{\raggedright
The inversion functions accept \code{Assumptions},
\code{"ObservablePower"}, \code{"Method"},
\code{"MaxMultiIndices"}, and \code{"MaxTotalDegree"}. Defaults are
\code{True}, $1$, \code{"Lagrange"}, $100000$, and $1000$, respectively.
The second implemented method is \code{"FixedPoint"}.\par}

\begin{lstlisting}
Clear[a, b, x, y];
r = RealInverseSeries[a x + b x^2, {x, 0}, y, 3,
      Assumptions -> a > 0 && Element[b, Reals]];
FullSimplify[r["Expression"], a > 0 && y > 0]
(* y/a - b y^2/a^3 *)

rSquared = RealInverseSeries[x + x^2, {x, 0}, y, 6,
                            "ObservablePower" -> 2];
\end{lstlisting}
These are fixed-parameter asymptotic statements. They do not assert
uniformity as, for example, a positive leading coefficient tends to zero
or an increment $\delta_i$ tends to zero. Such limits change the
normalization or the number of layers needed at a fixed cutoff.

To cross-check methods and the residual:
\begin{lstlisting}
rA = RealInverseSeries[x + x^2 (1 + Log[x]), {x, 0}, y, 5];
rB = RealInverseSeries[x + x^2 (1 + Log[x]), {x, 0}, y, 5,
                      "Method" -> "FixedPoint"];
FullSimplify[rA["Expression"] - rB["Expression"], y > 0]
(* 0 *)
RealInverseResidual[rA, 4]["AllZero"]
(* True: relative weights below 4 vanish *)
\end{lstlisting}
The comments above specify the mathematical expected outputs and are
covered by the supplied test suite; they are not presented as a transcript
from a native kernel run in this environment.

\subsection{Accepted syntax and explicit failure cases}

The expression parser accepts finite sums of products of constants,
powers $x^r$, and nonnegative integer powers of $\log x$. It expands sums
and products and then combines equal powers. A nonconstant polynomial in
$\log x$ is allowed in correction sectors but not as the leading
coefficient. The source variable and target variable must be distinct.

Unsupported inputs return a structured \code{Failure} through the
supported call signatures. Examples include a zero function, a
nonpositive leading power, an unproved positive leading coefficient,
complex coefficients, floating-point coefficients or exponents,
transcendental exponents such as $\pi$, nonpositive correction increments,
negative logarithmic powers, nested logarithms, and unknown method names.
A cutoff at or below the leading target power is rejected. Resource
limits stop the computation with a failure rather than return an
unannounced partial expansion.

For instance, $x+x^\pi$ is mathematically covered by the theorem but not
by the current exact-algebraic exponent backend. This is a deliberate
implementation boundary, not a claim that its inverse cannot be expanded:
\cref{eq:pure-product} applies directly with $\alpha=\pi$. Extending the
backend requires a trustworthy comparison and equality oracle for the
chosen larger class of constants.

\section{Implementation structure, invariants, and cost}
\label{sec:implementation}

\subsection{The construction pipeline}

The parser first produces a finite list of exponent--polynomial pairs.
It validates exact real algebraic exponents, sorts them by exact order,
combines collisions, and identifies the least exponent. That least
coefficient must be a provably positive constant. Dividing by it yields
the normalized correction polynomials.

The inversion engine then computes $B=pb-q$ and enumerates the bounded
simplex \cref{eq:enumeration-bound} recursively. It finds the least weight
at least $B$, keeps every multi-index at or below that frontier, computes
\cref{eq:master}, and merges equal weights. The fixed-point backend uses
the same cutoff and model but constructs the coefficients using only
finite jet multiplication, logarithm, binomial power, and composition.

The final stage simplifies coefficient polynomials under the supplied
assumptions, translates retained terms into $y$, and attaches the
first-omitted-layer metadata. Exact monomial inputs and $q=0$ return early.
The package never calls \code{Solve}, \code{InverseFunction},
\code{InverseSeries}, or \code{PowerExpand} to choose or construct the
inverse branch.

\subsection{Core invariants}

Every internal jet is a finite list of real algebraic weights with
polynomial amplitudes. Equal weights are merged. Before a logarithm or
binomial expansion, the nonconstant part has strictly positive weight.
All arithmetic is truncated only after the relevant weight has been
computed exactly. A source-level exponent must not survive validation as
an unresolved symbolic inequality.

\code{RootReduce} supplies canonical algebraic reduction for exponent
differences; its documented purpose and canonical behavior are described
in \cite{wl-rootreduce}. The comparator tests the sign of the reduced
exact difference. It never substitutes machine-precision approximations
as a tie-breaker. Polynomial coefficients are simplified separately, so
exponent arithmetic and coefficient algebra are not conflated.

Within the positive-power filtration, removing terms above a cutoff is
safe for subsequent multiplication and normalized composition: those
operations cannot decrease relative weight. Derivatives used in the
explicit coefficient theorem are applied inside the closed formula
before the resulting target power is assigned. Blindly differentiating a
truncated absolute-power jet and retaining its old cutoff would not be
safe; the package does not do that.

\subsection{Size bounds and practical limits}

If $K$ is the number of enumerated multi-indices and
$U=n_*\delta_*$, a simple bound is
\[
 K\le\prod_{i=1}^m\left(1+\left\lfloor\frac U{\delta_i}\right\rfloor\right).
\]
The simplex constraint is usually much smaller than this box bound.
Nevertheless, many sectors or a very small $\delta_*$ can make a modest
power cutoff expensive. Rational dependencies may cause many
multi-indices to merge into comparatively few final weights; the
multi-index engine still has to account for their contributions.

For a multi-index of degree $n$, the amplitude degree is at most
$\sum_i k_i d_i$. After forming its polynomial product, the $n-1$
first-order operators can be applied by linear passes over coefficient
vectors as in \cref{eq:operator-vector}. The reference implementation
favors transparent polynomial expressions over a specialized array
backend. Its jet multiplication uses direct convolution and repeated
sorting/merging; it is not an asymptotically optimized Newton-series
library.

The explicit method is normally the natural first choice. The fixed-point
method is valuable as an independent implementation and for studying
composition, but repeated polynomial convolutions can cost more. For
large problems, useful next steps are cached polynomial powers,
coefficient-vector arithmetic, a support-indexed sparse map, an
integer-lattice representation when a rational basis is available, and a
Newton backend. None of these optimizations changes the coefficient or
remainder theorem.

\section{Validation and reproduction}
\label{sec:validation}

\subsection{Executed checks}

The independent program \code{validation/validate.py} implements exact
weight comparison for its fixtures, the differential-operator formula,
and a separately coded truncated composition algebra. It uses SymPy
1.14.0 and mpmath 1.3.0 in the recorded run. Its 104 exact assertions
include the displayed logarithmic coefficients, the irrational-power
coefficients, agreement with the finite product formula, resonant
collisions, nonlinear leading powers, two distinct irrational
increments with mixed products, an irrational leading power, fractional
and negative observables, normalized composition residuals, and the
frontier sign-and-slope identity. Randomized fixtures use the fixed seed
$236367$ and are stored explicitly in \code{results.json}.

This program is independent of the Wolfram source. It is strong evidence
for the coefficient formulas and finite-jet algorithms and it generated
numerical examples reproducibly. It does \emph{not} execute Wolfram
syntax, test Wolfram evaluation rules, or validate the public parser.
Those are separate obligations of the MUnit suite.

For the numerical comparisons, the inverse was approximated by
180-decimal-digit bisection inside a positive bracket. The program
records the actual approximation error and its ratio to the first
omitted term. The following table uses the logarithmic expansion through
$y^4$, the irrational expansion through the $n=4$ correction, and
\cref{eq:nonlinear-core} through $t^2$. Values are rounded for display.
\begin{table}[htbp]
\centering\small
\begin{tabular}{@{}lrrr@{}}
\toprule
Family & $y$ & $|g-G|$ & $(g-G)/\text{frontier}$ \\
\midrule
Logarithmic & $10^{-2}$ & $1.48666\times10^{-7}$ & $1.0950093169$ \\
Logarithmic & $10^{-4}$ & $5.09642\times10^{-16}$ & $1.0022640153$ \\
Logarithmic & $10^{-8}$ & $1.16393\times10^{-34}$ & $1.0000005026$ \\
Logarithmic & $10^{-12}$ & $6.58956\times10^{-54}$ & $1.0000000001$ \\
Irrational power & $10^{-2}$ & $4.15553\times10^{-6}$ & $0.7805227225$ \\
Irrational power & $10^{-4}$ & $3.68347\times10^{-12}$ & $0.9597364481$ \\
Irrational power & $10^{-8}$ & $1.99266\times10^{-24}$ & $0.9990753832$ \\
Irrational power & $10^{-12}$ & $1.03646\times10^{-36}$ & $0.9999796056$ \\
Nonlinear core & $10^{-12}$ & $9.49699\times10^{-9}$ & $1.0887583784$ \\
Nonlinear core & $10^{-24}$ & $8.47300\times10^{-18}$ & $1.0017628576$ \\
Nonlinear core & $10^{-48}$ & $7.41352\times10^{-37}$ & $1.0000003656$ \\
\bottomrule
\end{tabular}
\caption{Executed 180-digit numerical comparisons. The finite approximations and first omitted terms are generated from exact coefficients.}
\label{tab:numerics}
\end{table}


The ratios approach $1$, as predicted by \cref{eq:first-omitted} when the
frontier polynomial is nonzero. At the less asymptotic test points they
need not be especially close to $1$: higher layers can still matter.
The exact proofs, not these samples, establish branch existence,
convergence, and the stated ordinary big-$O$ estimates.

\subsection{Native Wolfram execution status}

The remote Wolfram evaluator was attempted and returned
\code{invalid\_mcp\_response}, with an HTTP~404 response from its service
endpoint. A local Wolfram kernel was not available. Installing an
alternative interpreter also failed because network name resolution was
unavailable. Accordingly, no result in this deliverable is described as
a successful native Wolfram execution.

The Wolfram source received a string- and nested-comment-aware delimiter
scan and manual source review. These can detect some transcription
mistakes but are not substitutes for a Wolfram parser or evaluator.
The 95 MUnit tests in \code{Tests/RealInverseAsymptotics.wlt} are supplied
for native execution. They include public API expectations, structured
failure cases, resource limits, the two requested examples, and
cross-backend/residual/observable checks for the deterministic fixtures.
Their recorded status is \emph{not run}.

This distinction matters when assessing readiness. The mathematics has
self-contained proofs, and the underlying algorithms have executed
independent checks. The delivered Wolfram code remains a reference
implementation whose native integration tests must still be run in a
working kernel; it is not represented as production-certified software.

\subsection{Commands to reproduce the checks and the PDF}

From the extracted project directory:
\begin{lstlisting}
python validation/validate.py
wolframscript -file Tests/RunTests.wls
\end{lstlisting}
The Python command rewrites \code{validation/results.json} with the full
check list and numerical results. In a Wolfram notebook, the equivalent
native command is
\begin{lstlisting}
TestReport["path/to/RealInverseAsymptotics/Tests/RealInverseAsymptotics.wlt"]
\end{lstlisting}
The test file loads the package relative to its own location. The runner
prints succeeded and failed test counts and returns a nonzero exit code
unless the failed-test count is zero.

The article uses a standard pdf\LaTeX{} workflow, with Libertinus when
available and Latin Modern as a fallback:
\begin{lstlisting}
cd article
pdflatex -interaction=nonstopmode -halt-on-error real-inverse-asymptotics.tex
pdflatex -interaction=nonstopmode -halt-on-error real-inverse-asymptotics.tex
\end{lstlisting}
The source contains its bibliography directly; no external bibliography
database or proprietary font file is required. The numerical table is a
small included \TeX{} source derived from the recorded results. Font
files are not included in the archive.

\section{Extensions and boundaries of the present engine}
\label{sec:extensions}

\subsection{Other endpoints and asymptotic coordinates}

An inverse near a finite point $(x_0,y_0)$ can often be reduced to this
problem by writing $x=x_0+u$ and $y=y_0+v$ and choosing a positive side for
$u$ and $v$. The resulting leading term must fit $a u^p$ with $a,p>0$.
This is a caller-supplied coordinate change; the parser does not search
for centers or branches.

For an inverse at infinity, reciprocal variables can sometimes produce
an equivalent small-argument problem. One must transform the entire
function and its domain, not simply reuse an expansion with $y$ replaced
by $1/y$. Leading logarithmic or exponential cores commonly require a
different scale after such a transformation.

A left-hand real branch can also be reduced by $x=-u$ only when the
original powers and logarithms have been assigned consistent real
meanings on that side. For irrational powers, $x^\alpha$ on negative
$x$ is not automatically the same real-valued function as a signed power.
The package makes no such convention implicitly.

\subsection{Leading logarithms require a different core}

Consider
\[
 f(x)=x\log(1/x),\qquad x\to0^+.
\]
Here the logarithm belongs to the leading scale, not a positive-power
correction. With $u=-\log x$ the equation is $y=u\e^{-u}$, so the branch
approaching zero is
\begin{equation}
 x=-\frac{y}{W_{-1}(-y)}.
 \label{eq:leading-log-W}
\end{equation}
The branch $W_0$ would instead give $x$ near $1$. This is a particularly
clear example of why the branch must be fixed by its asymptotic
normalization rather than by the name of an inverse-function object.

For $a x^p(-\log x)^r$ with $r>0$, the exact core may similarly be written
\[
 u=-\frac rp W_{-1}\!\left(-\frac pr(y/a)^{1/r}\right),\qquad x=\e^{-u}.
\]
Subleading corrections can then be transported through this exact core
by phase-coordinate or scaled Lagrange inversion, as developed in the
supplied companion material \cite{companion}. The current package does
not implement this Lambert-core backend. It explicitly rejects a
logarithmic leading coefficient rather than return a power-core answer
outside its hypotheses.

\subsection{Infinite input expansions and flat corrections}

The finite analytic lifting is not, by itself, a convergence proof for
an arbitrary infinite forward expansion. For an infinite support, one
needs additional local-finiteness and summability conditions, or a
finite-jet statement with controlled remainders such as
\cref{eq:forward-remainder}. A list of plausible formal coefficients is
not a substitute for those conditions.

Flat corrections also contain information invisible to every finite
power--logarithmic jet. For example,
\[
 F(x)=a x^p(1+\e^{-1/x})
\]
has the same complete ordinary power--logarithmic asymptotic expansion
as $a x^p$, but its inverse satisfies
\[
 F^{-1}(y)=t-\frac tp\e^{-1/t}(1+o(1)),\qquad t=(y/a)^{1/p}.
\]
The coefficient follows from the leading displacement
$-\{F(t)-a t^p\}/(ap t^{p-1})$ and the smallness of the differentiated flat
term. The two inverses are not equal merely because all of their
power--logarithmic coefficients agree. Computing such a correction
requires keeping its exact sector in the input and using a larger
algebra; the present parser intentionally rejects it.

\subsection{Toward a broader engine}

The existing implementation already separates three concerns that should
remain separate in an extension: a branch/coordinate normalizer, an
ordered coefficient algebra, and a remainder or residual contract.
Supporting negative powers of logarithms would enlarge the coefficient
algebra from polynomials to a suitable asymptotic field. Supporting
leading logarithms would replace the power core by an exact Lambert
coordinate. Supporting exponential sectors would change the support
ordering and differentiation rules. Supporting symbolic exponents would
require conditional comparisons and possibly several parameter chambers.

None of those changes should be hidden inside a generic simplification
call. A useful generalized solver should return the selected branch,
the scale in which it expanded, and the assumptions under which its
ordering decisions hold. The finite engine here provides a fully
specified base case for that architecture, rather than a claim that the
remaining cases are automatic.

\section{Relation to the supplied research materials}
\label{sec:provenance}

The supplied \code{series-and-transseries} directory contains a much
broader program concerning compositional inversion and transseries.
For this report, the most directly relevant inspected source was the
companion volume \emph{Combinatorial Transseries and Their Inverses}
\cite{companion}. In particular, its sections labeled
\code{ct:sec:phase} and \code{ct:sec:scaled} distinguish an exactly
invertible core from a perturbation and derive finite Lagrange formulas.
Its discussion of marker convergence correctly emphasizes that a local
identity near zero markers does not alone justify setting every marker
to $1$.

The snapshot inspected for those sections had Git blob identifier
\begin{center}\small\ttfamily
05c56ce5ee60788a67f365110b62810550cba817.
\end{center}
The reference is to selected sections, not an assertion that the entire
large directory or all of its research-frontier claims were audited.

The construction in this article specializes that strategy to a
positive power core and polynomial logarithmic amplitudes, eliminates
the derivatives in the original variable through the operators
$q+\Delta+pj+\partial_\ell$, and supplies a finite analytic lifting and
an exact cutoff algorithm for this class. The broader transseries
literature \cite{edgar-beginners,edgar-grids} motivates the scale-based
viewpoint; the elementary local-finiteness and analytic arguments here
make the stated results independent of unresolved questions about more
general supports.

Classical Lagrange inversion is not claimed as new. The value of the
present development is the concrete specialization, the branch and
remainder contracts, the two independent finite algorithms, and a
reproducible implementation package for the exact examples in the
question.

\appendix
\section{Compact reference algorithms}
\label{app:algorithms}

The following pseudocode summarizes the explicit backend. Its input
amplitudes have already been combined at equal increments, zero
amplitudes removed, and all branch and exactness checks performed.

\begin{lstlisting}
input: a > 0, p > 0, corrections (delta_i, P_i(L)), observable q
       strict target-power cutoff b > q/p
B       := p*b - q
delta0  := minimum delta_i
upper   := ceiling(B/delta0)*delta0
indices := all k in N^m with sum(k_i*delta_i) <= upper
lambda  := minimum weight among indices with weight >= B
jet     := empty map from exact weights to polynomials

for each k with weight Delta <= lambda:
    n := sum(k_i)
    if n == 0:
        jet[0] += 1
    else:
        P := product(P_i(L)^k_i)
        for j = 1,...,n-1:
            P := (q + Delta + p*j)*P + derivative(P,L)
        jet[Delta] += (-1)^n*q*P/(p^n*product(k_i!))

combine and simplify each coefficient polynomial
retain Delta < B; substitute t=(y/a)^(1/p), L=Log[y/a]/p
return finite expression and the full polynomial jet[lambda]
\end{lstlisting}

A sparse normalized jet is an associative map from exact relative powers
to polynomials. Its composition backend is summarized below. Every
operation is truncated to the same relative bound.
\begin{lstlisting}
V := 1
repeat floor(bound/delta0)+1 times, or until stabilization:
    U := V - 1
    logV := sum_{n>=1} (-1)^(n+1)*U^n/n
    H := sum_i t^delta_i * (1+U)^delta_i * P_i(L+logV)
    V := (1+H)^(-1/p)
return V^q
\end{lstlisting}
All infinite sums displayed in this pseudocode terminate after
truncation because their arguments have positive relative order.
The actual source also performs structured validation, resource checks,
exact exponent comparison, and final coefficient simplification.

\section{A checklist for adding a new test family}

A new family should specify the positive domain and the exact leading
core first. It should then compare the explicit and fixed-point methods,
compose the retained inverse into the exact forward model, and check the
first residual coefficient against \cref{eq:frontier-residual}. At least
one test should place the cutoff exactly on a generated weight and one
should place it between two weights. A family with rational dependencies
should include a deliberate exponent collision; one with logarithms
should check both the polynomial coefficient and its logarithmic
remainder degree.

Numerical checks should use higher precision than the expected inverse
error, a bracket selecting the intended real branch, and several
geometrically separated target values. An observed small residual must
be interpreted through the appropriate slope. A zero formal frontier
must not be mislabeled a nonzero asymptotic leading error. These checks
exercise the mathematical contract rather than merely compare printed
expressions.

\section{Archive contents}

\begin{longtable}{@{}p{0.49\textwidth}p{0.47\textwidth}@{}}
\toprule
Path & Purpose \\
\midrule
\endhead
\code{article/real-inverse-asymptotics.tex} & Main article source. \\
\code{article/real-inverse-asymptotics.pdf} & Rendered article. \\
\code{article/numeric-table.tex} & Reproducible numerical table. \\
\code{Kernel/RealInverseAsymptotics.wl} & Wolfram Language reference package. \\
\code{Kernel/init.m} & Relative package loader. \\
\code{Examples/Examples.wl} & Runnable example expressions. \\
\code{Tests/RealInverseAsymptotics.wlt} & Native MUnit test specifications. \\
\code{Tests/RunTests.wls} & Native test runner. \\
\code{validation/validate.py} & Independent exact and numerical checks. \\
\code{validation/results.json} & Executed check list and recorded data. \\
\code{validation/REPORT.md} & Scope and status of validation. \\
\code{validation/static-check.py} & Delimiter scan, explicitly not an interpreter. \\
\code{README.md}, \code{LICENSE} & Usage, limitations, and license for original material. \\
\bottomrule
\end{longtable}

\begin{thebibliography}{99}

\bibitem{question}
Vladimir Reshetnikov,
\emph{How to get an asymptotic of the real-valued branch of the inverse
function?}, Mathematica Stack Exchange, question 236367 (2020),
user-supplied archive and online question, accessed 7 September 2026.
\url{https://mathematica.stackexchange.com/questions/236367/how-to-get-an-asymptotic-of-the-real-valued-branch-of-the-inverse-function}

\bibitem{dlmf}
NIST Digital Library of Mathematical Functions,
\S1.10(vii), \emph{Inverse Functions}, including the Lagrange inversion
formula. Accessed 7 September 2026.
\url{https://dlmf.nist.gov/1.10.vii}

\bibitem{companion}
Vladimir Reshetnikov,
\emph{Combinatorial Transseries and Their Inverses}, consolidated
companion volume in the ProveIt repository, September 2026;
selected sections on phase-coordinate inversion, scaled local reversion,
marker convergence, and coefficient calculus. Source snapshot identified
in \cref{sec:provenance}.
\url{https://github.com/VladimirReshetnikov/ProveIt/blob/main/Analysis/FabiusFunction/docs/semi-formalized-research-frontiers/drafts/series-and-transseries/Combinatorial_Transseries_Inverses/Combinatorial_Transseries_Inverses.tex}

\bibitem{edgar-beginners}
G. A. Edgar,
\emph{Transseries for Beginners}, Real Analysis Exchange
\textbf{35} (2010), 253--310; arXiv:0801.4877.
\url{https://arxiv.org/abs/0801.4877}

\bibitem{edgar-grids}
G. A. Edgar,
\emph{Transseries: Ratios, Grids, and Witnesses}, arXiv:0909.2430 (2009).
\url{https://arxiv.org/abs/0909.2430}

\bibitem{wl-inverse}
Wolfram Research,
\emph{InverseSeries}, Wolfram Language \& System Documentation Center.
Accessed 7 September 2026.
\url{https://reference.wolfram.com/language/ref/InverseSeries.html}

\bibitem{wl-asymptotic}
Wolfram Research,
\emph{AsymptoticSolve}, Wolfram Language \& System Documentation Center.
Accessed 7 September 2026.
\url{https://reference.wolfram.com/language/ref/AsymptoticSolve.html}

\bibitem{wl-rootreduce}
Wolfram Research,
\emph{RootReduce}, Wolfram Language \& System Documentation Center.
Accessed 7 September 2026.
\url{https://reference.wolfram.com/language/ref/RootReduce.html}

\end{thebibliography}
\end{document}
